\chapter{Sự tiến hóa của trí tuệ nhân tạo: Từ mùa đông đến phục hưng}

\begin{quote}\small

\textbf{Tóm tắt.} Chương 1 phác họa quỹ đạo lịch sử của trí tuệ nhân tạo, từ những khởi đầu mang tính biểu tượng và sự lạc quan ban đầu, qua nỗi thất vọng của Mùa đông AI, đến thời kỳ phục hưng hiện nay được thúc đẩy bởi học sâu và các mô hình dựa trên transformer. Chương này xem xét các yếu tố đã tạo điều kiện cho sự trỗi dậy trở lại của AI, những tiến bộ về sức mạnh tính toán, khả năng sẵn có của dữ liệu và đổi mới thuật toán, và đạt đến đỉnh điểm với sự xuất hiện của AI tạo sinh như một bước phát triển mang tính thay đổi mô hình. Chương cũng đánh giá một cách phê phán các thách thức đang diễn ra, bao gồm chi phí môi trường, những hạn chế về hạ tầng, sự khan hiếm dữ liệu và các mối quan ngại về công bằng. Bằng cách đặt các đột phá đương đại trong bối cảnh lịch sử, chương nhấn mạnh sự cần thiết của đổi mới cân bằng và có trách nhiệm, tránh việc hứa hẹn quá mức từng làm đình trệ tiến bộ của lĩnh vực này.

\subsection{Từ khóa}

Mùa đông AI, học sâu, mạng nơ-ron, AI tạo sinh, mô hình transformer, sức mạnh tính toán, AI dựa trên dữ liệu, đổi mới có trách nhiệm

\section{}

\section{\texorpdfstring{\textbf{Mùa đông AI} \textbf{- Một góc nhìn lịch sử}}{Mùa đông AI - Một góc nhìn lịch sử}}

Câu chuyện về trí tuệ nhân tạo là câu chuyện của những đỉnh cao và vực sâu đầy kịch tính, và có lẽ không vực sâu nào sâu hơn giai đoạn được gọi là ``Mùa đông AI'' (Hendler, 2008). Sự ra đời của lĩnh vực này có thể được truy về một thời điểm mang tính bước ngoặt vào năm 1956 tại Dự án Nghiên cứu Mùa hè Dartmouth, nơi John McCarthy đặt ra thuật ngữ ``trí tuệ nhân tạo'' (McCarthy et al., 1955). Cuộc tụ họp của những bộ óc xuất chúng này, trong đó có Marvin Minsky, Claude Shannon và Herbert Simon, đã đánh dấu sự khởi đầu của AI với tư cách một ngành học chính thức.

Nền tảng cho lĩnh vực mới này đã được đặt ra trước đó bởi bài báo có tầm ảnh hưởng của Alan Turing, ``Computing Machinery and Intelligence'' (1950), công trình giới thiệu Phép thử Turing nổi tiếng và thiết lập những khái niệm căn bản về trí tuệ máy móc. Những nhà tiên phong ban đầu này hình dung ra những cỗ máy có thể tái tạo năng lực nhận thức của con người, từ giải quyết vấn đề đến hiểu ngôn ngữ tự nhiên (McCorduck, 2004).

Thập niên 1950 và đầu thập niên 1960 chứng kiến những thành công ban đầu đáng kể. Logic Theorist (1956) của Newell và Simon trở thành chương trình đầu tiên chứng minh được các định lý toán học, trong khi General Problem Solver (GPS) của họ thể hiện năng lực suy luận cơ bản (Newell et al., 1959). Những thành tựu này, cùng với các tiến bộ trong xử lý ngôn ngữ tự nhiên (NLP) như ELIZA (1966) của Weizenbaum, đã tạo nên một làn sóng lạc quan. Giai đoạn này nhận được nguồn tài trợ đáng kể từ cả các cơ quan chính phủ, đặc biệt là Văn phòng Kỹ thuật Xử lý Thông tin (IPTO) của DARPA, lẫn các phòng thí nghiệm nghiên cứu doanh nghiệp như Phòng thí nghiệm AI của MIT và Phòng thí nghiệm AI của Stanford (Roland \& Shiman, 2002).

Tuy nhiên, những mục tiêu đầy tham vọng của lĩnh vực này sớm vấp phải các trở ngại đáng kể. Những hạn chế của AI biểu tượng (symbolic AI), mô hình thống trị thời bấy giờ, ngày càng lộ rõ. Các hệ thống này, dựa trên logic hình thức và suy luận theo quy tắc, gặp khó khăn trước sự phức tạp của thế giới thực, đặc biệt trong những lĩnh vực đòi hỏi suy luận thường thức hoặc nhận dạng mẫu (Dreyfus, 1972). Cách tiếp cận ``thế giới vi mô'' (microworld), tiêu biểu là chương trình SHRDLU của Winograd, tuy thành công trong các môi trường bị giới hạn nhưng đã không thể mở rộng sang những kịch bản thực tế phức tạp hơn (Winograd, 1972).

Đến giữa thập niên 1970, làn sóng chỉ trích dồn dập lên đến đỉnh điểm qua một số báo cáo mang tính tàn phá. Báo cáo Lighthill (1973), do Hội đồng Nghiên cứu Khoa học Anh ủy nhiệm, đặc biệt chỉ trích việc AI không đạt được những ``mục tiêu hoành tráng'' của mình và đặt nghi vấn về giá trị thực tiễn của lĩnh vực này. Tác động của báo cáo diễn ra ngay lập tức và nghiêm trọng, dẫn đến việc kinh phí nghiên cứu AI ở Anh gần như bị xóa bỏ hoàn toàn (Lighthill, 1973). Sự hoài nghi tương tự lan rộng khắp Hoa Kỳ, khi DARPA và các nhà tài trợ lớn khác ngày càng vỡ mộng về tiến độ của AI (Crevier, 1993).

Mùa đông AI tiếp theo của thập niên 1970 và 1980 được đặc trưng bởi một chuỗi thất bại dồn dập. Sự sụp đổ của thị trường thương mại dành cho các hệ thống AI, đặc biệt là hệ chuyên gia, đã dẫn đến việc giải thể nhiều công ty và chương trình nghiên cứu AI. XCON, do Digital Equipment Corporation phát triển, minh họa cho những thách thức này: dù ban đầu thành công, chi phí bảo trì của nó tăng vọt khi hệ thống ngày càng phức tạp (McDermott, 1982). Những hạn chế của hệ chuyên gia trở nên rõ ràng: chúng dễ vỡ, không thể học từ kinh nghiệm và đòi hỏi quá trình kỹ nghệ tri thức thủ công tốn nhiều công sức.

Giai đoạn thu hẹp này, tuy đau đớn, đã tỏ ra vô giá đối với sự phát triển của lĩnh vực. Nó thúc đẩy một cách tiếp cận nghiêm ngặt hơn đối với nghiên cứu AI, nhấn mạnh việc đánh giá thực nghiệm và đặt ra các mục tiêu thực tế. Những năm mùa đông cũng chứng kiến sự xuất hiện âm thầm của các phương pháp mà sau này sẽ làm thay đổi căn bản lĩnh vực, đặc biệt trong mạng nơ-ron và suy luận xác suất (Nilsson, 2009). Máy kết nối (connection machine), do Danny Hillis phát triển vào thập niên 1980, đại diện cho những nỗ lực ban đầu về kiến trúc xử lý song song, thứ về sau tỏ ra then chốt đối với học sâu (Hillis, 1985).

Mùa đông AI đã dạy cho lĩnh vực này những bài học quan trọng về mối nguy của việc hứa hẹn quá mức và tầm quan trọng của việc quản lý kỳ vọng. Nó cho thấy rằng tiến bộ trong AI đòi hỏi không chỉ những đổi mới về thuật toán mà còn cả những bước tiến trong hạ tầng tính toán và hiểu biết lý thuyết. Có lẽ quan trọng nhất, nó nhấn mạnh sự cần thiết của một sự đánh giá tinh tế hơn về bản thân trí tuệ, tức là sự nhận ra rằng nhận thức của con người phức tạp hơn nhiều so với giả định ban đầu, đòi hỏi nhiều cách tiếp cận và tiến bộ bền bỉ, tăng dần thay vì những đột phá nhanh chóng (Brooks, 1991).

\section{Sự hồi sinh của AI: Các yếu tố tạo nên làn sóng bùng nổ hiện nay}

Giai đoạn ``tan băng'' của lĩnh vực AI bắt đầu một cách dần dần vào cuối những năm 1990 và đầu những năm 2000, nhờ một số tiến bộ đồng thời về công nghệ và hạ tầng. Bước đột phá quan trọng nhất là sự hồi sinh của mạng nơ-ron thông qua học sâu, một sự hồi sinh sẽ làm thay đổi căn bản lĩnh vực này. Thời kỳ phục hưng ấy được thúc đẩy bởi ba yếu tố then chốt hội tụ lại, tạo nên ``cơn bão hoàn hảo'' cho sự phát triển của AI: sức mạnh tính toán chưa từng có, nguồn dữ liệu khổng lồ và những đổi mới mang tính cách mạng về thuật toán (Goodfellow et al., 2016).

\subsection{Cuộc cách mạng tính toán: GPU và hơn thế nữa}

Sự trỗi dậy của Bộ xử lý đồ họa (Graphics Processing Units, GPU) đánh dấu một thời điểm mang tính chuyển đổi trong bức tranh tính toán. Ban đầu được thiết kế để kết xuất đồ họa 3D phức tạp trong trò chơi điện tử, các bộ xử lý này sở hữu một kiến trúc độc đáo mà tình cờ lại đặc biệt lý tưởng cho nhu cầu xử lý song song của mạng nơ-ron. Khác với các Bộ xử lý trung tâm (Central Processing Units, CPU) truyền thống vốn được thiết kế để xử lý tuần tự, GPU có thể thực hiện hàng nghìn phép tính đồng thời, một khả năng hoàn toàn phù hợp với các phép toán ma trận vốn là nền tảng của việc huấn luyện mạng nơ-ron (Oh \& Jung, 2004).

Việc NVIDIA giới thiệu CUDA (Compute Unified Device Architecture) vào năm 2006 là một thời điểm then chốt, cung cấp cho các nhà phát triển quyền truy cập trực tiếp vào năng lực tính toán song song của GPU. Bước phát triển này đã biến GPU từ phần cứng chuyên dụng cho trò chơi thành những cỗ máy tính toán đa dụng mạnh mẽ. Sau đó, công ty nhận ra tiềm năng đang nổi lên của các ứng dụng AI và bắt đầu phát triển những kiến trúc chuyên biệt. Việc ra mắt kiến trúc Tesla vào năm 2007, tiếp theo là các thế hệ phần cứng được tối ưu hóa cho AI, đã rút ngắn đáng kể thời gian huấn luyện và cho phép phát triển các mô hình ngày càng phức tạp (Kirk \& Hwu, 2017).

Tác động là sâu sắc: những tác vụ từng đòi hỏi hàng tuần tính toán nay có thể hoàn thành trong vài ngày hoặc vài giờ. Chẳng hạn, việc huấn luyện AlexNet vào năm 2012, một cột mốc mang tính bước ngoặt của học sâu, đã tận dụng hai GPU NVIDIA GTX 580 để đạt hiệu năng đột phá trong nhận dạng hình ảnh, rút ngắn thời gian huấn luyện từ vài tuần xuống còn vài ngày (Krizhevsky et al., 2012).

\subsection{Cuộc cách mạng dữ liệu: Từ khan hiếm đến dồi dào}

Cuộc cách mạng số mang đến một thành tố then chốt khác: lượng dữ liệu dồi dào chưa từng có. Sự tăng trưởng bùng nổ của internet, sự lan rộng của các nền tảng mạng xã hội và quá trình số hóa trên diện rộng các ngành công nghiệp đã tạo ra những tập dữ liệu khổng lồ, phù hợp để huấn luyện các hệ thống AI. Cuộc cách mạng dữ liệu này được đặc trưng bởi những bước phát triển mang tính biến đổi trên nhiều lĩnh vực. Sự xuất hiện của các nền tảng mạng xã hội đã tạo ra lượng nội dung do người dùng tạo lập khổng lồ, cung cấp dữ liệu ngôn ngữ tự nhiên ở quy mô chưa từng có. Song song với đó, việc áp dụng rộng rãi các thiết bị di động trang bị cảm biến tinh vi đã tạo ra những tập dữ liệu phong phú về hình ảnh, dữ liệu vị trí và các tương tác của người dùng. Internet vạn vật (IoT) còn non trẻ bắt đầu tạo ra các luồng dữ liệu cảm biến liên tục từ hàng triệu thiết bị được kết nối, trong khi các sáng kiến chuyển đổi số trong nhiều ngành đã sinh ra các tập dữ liệu có cấu trúc từ y tế, tài chính và các quy trình công nghiệp.

Các nền tảng điện toán đám mây nổi lên như một lực lượng dân chủ hóa, thay đổi căn bản cách các nhà nghiên cứu và tổ chức có thể truy cập và xử lý những tập dữ liệu khổng lồ này. Các dịch vụ như Amazon Web Services (AWS), Google Cloud Platform và Microsoft Azure đã loại bỏ nhu cầu đầu tư hạ tầng ban đầu quy mô lớn. Sự dân chủ hóa này giúp các công ty khởi nghiệp và các nhà nghiên cứu có thể thử nghiệm AI ở những quy mô trước đây chỉ dành cho các công ty công nghệ lớn (Jonas et al., 2019).

\subsection{Những đột phá về thuật toán: Từ lý thuyết đến thực tiễn}

Sự hội tụ của các lợi thế về tính toán và dữ liệu đã kết hợp một cách cộng hưởng với những đột phá lý thuyết trong các thuật toán học máy. Trong lĩnh vực đổi mới kiến trúc, việc hoàn thiện các kỹ thuật lan truyền ngược (backpropagation) đã giúp việc huấn luyện các mạng nơ-ron sâu trở nên khả thi hơn, trong khi sự ra đời của mạng nơ-ron tích chập (CNN) đã tạo nên cuộc cách mạng cho các tác vụ thị giác máy tính. Đồng thời, mạng nơ-ron hồi quy (RNN) và mạng bộ nhớ dài-ngắn hạn (LSTM) cho phép xử lý dữ liệu tuần tự tốt hơn, được bổ sung bởi các kỹ thuật dropout và chuẩn hóa theo lô (batch normalization) vốn giải quyết những thách thức then chốt trong huấn luyện.

Những tiến bộ đáng kể trong các kỹ thuật tối ưu hóa càng thúc đẩy sự phát triển của lĩnh vực này. Việc phát triển các thuật toán tối ưu hóa cải tiến như Adam giúp huấn luyện nhanh hơn và ổn định hơn, trong khi các kỹ thuật khởi tạo tốt hơn góp phần giải quyết vấn đề triệt tiêu gradient (vanishing gradient). Việc giới thiệu các hàm kích hoạt mới, đặc biệt là Rectified Linear Unit (ReLU), cho phép huấn luyện nhanh hơn và đạt hiệu suất tốt hơn trong nhiều ứng dụng đa dạng.

Hiểu biết về động lực học của quá trình huấn luyện cũng được cải thiện đáng kể trong giai đoạn này. Các nhà nghiên cứu có được những hiểu biết sâu sắc hơn về địa hình hàm mất mát (loss landscape), từ đó dẫn đến các kiến trúc mạng được tối ưu hóa. Các kỹ thuật điều chuẩn (regularization) tiên tiến làm giảm hiện tượng quá khớp (overfitting), trong khi sự phát triển của các phương pháp học chuyển giao (transfer learning) cho phép sử dụng các mô hình được huấn luyện trước một cách hiệu quả hơn. Sự hiểu biết lý thuyết được nâng cao này tạo điều kiện cho việc xây dựng các kiến trúc mạng nơ-ron mạnh mẽ và hiệu quả hơn.

Sự cộng hưởng giữa ba trụ cột này, tính toán, dữ liệu và thuật toán, đã tạo ra một vòng phản hồi mạnh mẽ. Phần cứng mạnh hơn cho phép xây dựng các mô hình và tập dữ liệu lớn hơn, điều này thúc đẩy các cải tiến thuật toán, và đến lượt mình lại kích thích việc tối ưu hóa phần cứng. Vòng tuần hoàn tích cực này tiếp tục đẩy nhanh sự phát triển của AI cho đến ngày nay, dẫn đến các mô hình và ứng dụng ngày càng tinh vi (LeCun et al., 2015).

\subsection{Sự trỗi dậy của AI tạo sinh: Một sự chuyển dịch mô hình}

Sự xuất hiện của AI tạo sinh đánh dấu một chương mang tính biến đổi trong lịch sử trí tuệ nhân tạo, được xây dựng trên nhiều đổi mới đột phá về kiến trúc học sâu và phương pháp huấn luyện. Sự chuyển dịch này bắt đầu với việc giới thiệu kiến trúc transformer vào năm 2017, vốn đã cách mạng hóa một cách căn bản cách các hệ thống AI xử lý dữ liệu tuần tự. Đổi mới then chốt là cơ chế chú ý (attention), cho phép mô hình cân nhắc động tầm quan trọng của các phần tử đầu vào khác nhau, đã tỏ ra hiệu quả đáng kể trên nhiều lĩnh vực, vượt qua các phương pháp trước đó cả về hiệu suất lẫn năng lực (Vaswani et al., 2017).

Kiến trúc transformer đã đưa ra nhiều đổi mới quan trọng trong thiết kế mạng nơ-ron. Cơ chế tự chú ý (self-attention) cho phép mô hình hóa trực tiếp mối quan hệ giữa mọi phần tử trong một chuỗi, bất kể khoảng cách giữa chúng, qua đó giải quyết hiệu quả các vấn đề phụ thuộc tầm xa vốn đã gây khó khăn cho các kiến trúc trước đây. Cơ chế chú ý đa đầu (multi-head attention) tăng cường thêm năng lực này bằng cách cho phép mô hình nắm bắt đồng thời nhiều kiểu quan hệ khác nhau. Những đổi mới kiến trúc này được bổ sung bởi mã hóa vị trí (positional encoding), giúp bảo toàn thông tin về thứ tự của chuỗi đồng thời cho phép xử lý song song (Devlin et al., 2019).

Các mô hình ngôn ngữ lớn (Large Language Models, LLM) nổi lên như câu chuyện thành công lớn đầu tiên của các kiến trúc dựa trên transformer, thể hiện những năng lực chưa từng có trong xử lý ngôn ngữ tự nhiên (NLP). Quá trình tiến triển từ GPT gốc (Generative Pre-trained Transformer) qua các thế hệ kế tiếp đã minh họa khả năng mở rộng đáng kinh ngạc của kiến trúc transformer. Mỗi phiên bản đều mang lại những cải thiện đáng kể về khả năng hiểu và tạo sinh ngôn ngữ, với các mô hình thể hiện những hành vi ngày càng tinh vi như học trong ngữ cảnh (in-context learning) và khái quát hóa tác vụ. Các quy luật mở rộng (scaling laws) được khám phá trong giai đoạn này cho thấy một mối quan hệ có hệ thống giữa kích thước mô hình, khối lượng dữ liệu huấn luyện và hiệu năng, gợi ý rằng các mô hình lớn hơn, được huấn luyện trên các tập dữ liệu rộng hơn, có thể đạt được những năng lực khác biệt về chất (Kaplan et al., 2020).

Sự phát triển của BERT (Bidirectional Encoder Representations from Transformers) đã đưa vào khả năng hiểu ngữ cảnh hai chiều, cho phép thấu hiểu sắc thái hơn về cấu trúc và ý nghĩa của ngôn ngữ. Đổi mới này làm nảy sinh nhiều biến thể kiến trúc được tối ưu hóa cho các ứng dụng cụ thể, từ việc tối ưu hóa mạnh mẽ của RoBERTa đến thiết kế tiết kiệm tham số của ALBERT (Liu et al., 2019). Các mô hình này thể hiện năng lực đáng kể trong những tác vụ từ trả lời câu hỏi đến phân loại văn bản, đồng thời cũng hé lộ những hiểu biết về cách mạng nơ-ron xử lý và biểu diễn ngôn ngữ.

Học chuyển giao (transfer learning) nổi lên như một kỹ thuật then chốt trong việc triển khai các mô hình lớn này, làm thay đổi căn bản cách các hệ thống AI được phát triển và triển khai. Cách tiếp cận này cho phép các mô hình được tiền huấn luyện trên những tập dữ liệu khổng lồ được tinh chỉnh cho các tác vụ cụ thể chỉ với lượng dữ liệu bổ sung tương đối ít, giảm đáng kể yêu cầu về tính toán và dữ liệu cho các ứng dụng mới. Sự xuất hiện của năng lực học vài mẫu (few-shot) và học không mẫu (zero-shot) càng chứng tỏ khả năng của các mô hình này trong việc thích ứng với tác vụ mới với rất ít hoặc không cần hướng dẫn tường minh, gợi ý một dạng siêu học tập (meta-learning) gần giống hơn với sự linh hoạt nhận thức của con người (Brown et al., 2020).

Sau đó, lĩnh vực này mở rộng ra ngoài văn bản để đón nhận các mô hình đa phương thức có khả năng xử lý và tạo sinh trên nhiều dạng phương tiện khác nhau. Sự tiến hóa này được đánh dấu bằng việc phát triển các kiến trúc có thể tích hợp liền mạch nhiều loại dữ liệu, từ hình ảnh và âm thanh đến video và thông tin có cấu trúc. DALL-E đã chứng minh khả năng tạo ra hình ảnh từ mô tả bằng văn bản, trong khi các mô hình như GPT-4V cho thấy năng lực hiểu và suy luận về thông tin thị giác song song với văn bản. Stable Diffusion đã giới thiệu các kỹ thuật mới để tạo ảnh hiệu quả, giúp các nhà nghiên cứu và nhà phát triển dễ tiếp cận hơn những năng lực này (Ramesh et al., 2022).

Tác động của sự mở rộng đa phương thức này là sâu sắc, mở ra các ứng dụng mới trong công cụ sáng tạo, trực quan hóa khoa học và tương tác người-máy. Những mô hình này thể hiện khả năng hiểu và tạo sinh xuyên các phương thức theo những cách gần với năng lực nhận thức của con người hơn. Sự phát triển của các kỹ thuật tạo sinh có kiểm soát cũng cho phép đầu ra chính xác và đáng tin cậy hơn, giải quyết những thách thức trước đây về tính nhất quán và tính mạch lạc.

Các đổi mới kiến trúc gần đây tập trung vào việc cải thiện hiệu suất và khả năng kiểm soát trong khi vẫn duy trì chất lượng. Việc phát triển các cơ chế chú ý thưa (sparse attention), như những cơ chế được dùng trong Longformer và BigBird, đã cho phép xử lý các chuỗi dài hơn nhiều đồng thời giảm yêu cầu tính toán. Trong khi đó, các kỹ thuật kiểm soát và định hướng quá trình tạo sinh ngày càng tinh vi, cho phép đầu ra đáng tin cậy và có chủ đích hơn (Zaheer et al., 2020).

Những tiến bộ này là các bước tiến đáng kể hướng tới trí tuệ nhân tạo tổng quát hơn, mặc dù những thách thức quan trọng vẫn còn. Các vấn đề về thiên kiến, độ tin cậy và hiệu quả tính toán tiếp tục thúc đẩy nghiên cứu về các kiến trúc và phương pháp huấn luyện tinh vi hơn. Lĩnh vực này cũng bắt đầu phải đối mặt với các câu hỏi về khả năng diễn giải và kiểm soát mô hình, khi những hệ thống này ngày càng có năng lực và được triển khai rộng rãi.

\section{Những thách thức hiện tại và mối quan ngại trong tương lai}

Bất chấp những tiến bộ vượt bậc của trí tuệ nhân tạo, lĩnh vực này vẫn đối mặt với những thách thức đáng kể, đe dọa sự phát triển bền vững và việc triển khai công bằng của nó. Những thách thức này trải rộng trên các phương diện môi trường, kinh tế, kỹ thuật và xã hội, đòi hỏi sự cân nhắc cẩn trọng và các giải pháp mang tính hệ thống.

Tác động môi trường của việc huấn luyện các mô hình AI lớn đã nổi lên như một mối quan ngại cấp bách trong cộng đồng khoa học. Các nghiên cứu đương đại cho thấy lượng khí thải carbon đáng kể gắn với việc phát triển và triển khai các hệ thống AI tiên tiến. Một lần huấn luyện duy nhất cho một mô hình ngôn ngữ lớn (LLM) có thể tạo ra lượng khí thải carbon tương đương với lượng khí thải hằng năm của hàng trăm chiếc ô tô, và một số ước tính cho rằng lượng phát thải có thể lên tới 626.000 pound carbon dioxide tương đương (Strubell et al., 2019). Chi phí môi trường này bắt nguồn từ nguồn lực tính toán khổng lồ cần thiết cho việc huấn luyện mô hình, vốn thường bao gồm nhiều vòng lặp thử nghiệm và các kiến trúc ngày càng lớn hơn.

Hạ tầng phần cứng hỗ trợ phát triển AI là một thách thức then chốt khác. Tình trạng thiếu hụt chất bán dẫn toàn cầu đã phơi bày những điểm yếu trong chuỗi cung ứng phần cứng tính toán chuyên dụng cho AI. Sự khan hiếm này ảnh hưởng đến mọi thứ, từ bộ xử lý đồ họa (GPU) đến các bộ tăng tốc AI chuyên biệt, dẫn đến chi phí tăng và tiến độ nghiên cứu bị chậm trễ. Những hạn chế về sản xuất, đặc biệt ở các tiến trình bán dẫn tiên tiến, đã tạo ra các nút thắt cổ chai trong việc sản xuất các chip tối ưu hóa cho AI. Hơn nữa, căng thẳng địa chính trị xoay quanh sản xuất chất bán dẫn và kiểm soát xuất khẩu có nguy cơ làm phân mảnh cộng đồng nghiên cứu AI toàn cầu và cản trở sự phát triển mang tính hợp tác (Khan, 2022).

Sự khan hiếm và chất lượng của dữ liệu đặt ra những trở ngại ngày càng đáng kể khi lĩnh vực này tiến triển. Trong khi giai đoạn đầu phát triển AI được hưởng lợi từ các bộ dữ liệu sẵn có, thì cuộc đua giành dữ liệu huấn luyện chất lượng cao đã trở nên ngày càng cạnh tranh. Các nguồn dữ liệu trước đây có thể tiếp cận đang dần cạn kiệt, đặc biệt trong các lĩnh vực chuyên biệt và các ngôn ngữ có sự hiện diện số hạn chế. Sự khan hiếm này tác động không cân xứng đến nghiên cứu về các ngôn ngữ ít tài nguyên và các lĩnh vực khoa học chuyên ngành, nơi các bộ dữ liệu được chú thích đặc biệt hiếm. Vấn đề không chỉ dừng ở số lượng mà còn mở rộng sang các câu hỏi về chất lượng dữ liệu, tính đại diện và các cân nhắc đạo đức trong thu thập dữ liệu (Sambasivan et al., 2021).

Phương diện kinh tế của phát triển AI đã tạo ra những mô hình tập trung đáng lo ngại. Nguồn lực tính toán và tài chính cần thiết để huấn luyện và triển khai các mô hình AI hiện đại nhất đã tăng vọt, với các ước tính cho rằng huấn luyện một mô hình ngôn ngữ lớn duy nhất có thể tốn hàng triệu đô la. Rào cản kinh tế này đã dẫn đến việc năng lực AI tập trung trong tay một số ít tổ chức có nhiều nguồn lực, chủ yếu là các công ty công nghệ lớn và các viện nghiên cứu tinh hoa. Sự tập trung này làm dấy lên những lo ngại về công bằng trong tiếp cận và tính đa dạng của các quan điểm trong phát triển AI (Ahmed \& Wahed, 2020).

Các thách thức kỹ thuật vẫn tồn tại trong việc mở rộng quy mô hệ thống AI một cách hiệu quả. Dù các mô hình lớn hơn đã thể hiện những năng lực ấn tượng, chúng cũng bộc lộ những hạn chế về khả năng suy luận, độ tin cậy và tính vững chắc. Hiện tượng mô hình dễ vỡ, khi các hệ thống tưởng chừng có năng lực lại thất bại bất ngờ trước các trường hợp biên, vẫn là một mối quan ngại đáng kể. Ngoài ra, các kiến trúc hiện nay gặp khó khăn với những vấn đề đòi hỏi suy luận theo lẽ thường hoặc hiểu biết nhân quả, bất chấp hiệu suất mạnh mẽ của chúng trong các tác vụ nhận dạng mẫu (Pearl \& Mackenzie, 2018).

Những hạn chế về hạ tầng là một thách thức đáng kể khác. Kích thước ngày càng tăng của các mô hình AI gây áp lực lên hạ tầng tính toán, băng thông mạng và hệ thống lưu trữ hiện có. Nhu cầu huấn luyện phân tán giữa các trung tâm dữ liệu đòi hỏi các hệ thống điều phối tinh vi và mạng băng thông cao, trong khi việc triển khai mô hình đối mặt với các thách thức về độ trễ và hiệu quả sử dụng tài nguyên. Những yêu cầu hạ tầng này càng làm trầm trọng thêm sự tập trung năng lực AI vào các tổ chức có nhiều nguồn lực.

Các mối quan ngại về quyền riêng tư và an ninh cũng đã nổi lên như những thách thức quan trọng. Nhu cầu thu thập và chia sẻ dữ liệu quy mô lớn phải được cân bằng với quyền riêng tư của cá nhân và các quy định về bảo vệ dữ liệu. Ngoài ra, khả năng các hệ thống AI bị xâm phạm hoặc bị lạm dụng làm dấy lên những lo ngại đáng kể về an ninh, đặc biệt khi các hệ thống này được triển khai trong các ứng dụng nhạy cảm (Papernot et al., 2018).

Thách thức về khả năng tái lập trong nghiên cứu AI ngày càng trở nên rõ rệt. Sự kết hợp giữa yêu cầu tính toán quy mô lớn, các bộ dữ liệu độc quyền và quy trình huấn luyện phức tạp khiến các nhà nghiên cứu khó kiểm chứng và phát triển tiếp từ các kết quả đã công bố. Cuộc khủng hoảng tái lập này đe dọa nền tảng khoa học của lĩnh vực và cản trở tiến bộ trong việc giải quyết các thách thức kỹ thuật khác (Pineau et al., 2021).

\section{Con đường phía trước: Cân bằng giữa đổi mới và trách nhiệm}

Tương lai của việc phát triển AI phải tìm được sự cân bằng giữa đổi mới liên tục và việc quản lý có trách nhiệm các nguồn lực. Nghiên cứu về các phương pháp huấn luyện hiệu quả hơn, chẳng hạn các cơ chế chú ý thưa (sparse attention) và những mô hình nhỏ hơn, chuyên biệt hơn, mở ra triển vọng giảm nhu cầu tính toán. Các đổi mới phần cứng trong điện toán mô phỏng thần kinh (neuromorphic computing) và bộ xử lý lượng tử có thể mang lại những hướng đi mới cho các hệ thống AI tiết kiệm năng lượng.

Dân chủ hóa khả năng tiếp cận AI đòi hỏi phải giải quyết cả các rào cản kỹ thuật lẫn kinh tế. Các sáng kiến mã nguồn mở, các phương pháp học liên kết (federated learning) và các kỹ thuật nén mô hình được cải tiến có thể giúp phân bổ năng lực AI rộng rãi hơn. Các khuôn khổ hợp tác quốc tế nhằm chia sẻ nghiên cứu, dữ liệu và tài nguyên tính toán có thể góp phần ngăn chặn sự tập trung nguy hiểm của năng lực AI.

Lĩnh vực này cũng phải chấp nhận các thực hành đổi mới có trách nhiệm. Điều đó bao gồm việc xây dựng các biện pháp an toàn vững chắc, bảo đảm các hệ thống AI tôn trọng quyền riêng tư và nhân quyền, đồng thời thiết lập các khuôn khổ quản trị thúc đẩy sự phát triển AI có lợi trong khi kiểm soát rủi ro. Những bài học của Mùa đông AI nhắc chúng ta rằng tiến bộ bền vững đòi hỏi phải cân bằng giữa tham vọng và tính thực tế, bảo đảm rằng việc phát triển AI phục vụ lợi ích rộng lớn hơn của nhân loại thay vì chỉ những thành tựu kỹ thuật hẹp hòi.

Khi tiến về phía trước, thách thức then chốt sẽ là duy trì đà phát triển hiện tại trong khi giải quyết những thách thức nền tảng này. Thành công sẽ đòi hỏi sự hợp tác chưa từng có giữa các nhà nghiên cứu, lãnh đạo ngành, nhà hoạch định chính sách và công chúng nói chung, nhằm bảo đảm rằng tiềm năng của AI được hiện thực hóa theo cách mang lại lợi ích cho toàn thể nhân loại.

\section{Kết luận}

Lịch sử của trí tuệ nhân tạo phản ánh một chu kỳ lặp đi lặp lại: những tham vọng lớn lao được theo sau bởi các giai đoạn vỡ mộng. Mùa đông AI (AI Winter) đã để lại những bài học quan trọng về việc quản lý kỳ vọng, nhấn mạnh tầm quan trọng của các mục tiêu thực tế và tiến bộ từng bước. Sự phục hưng gần đây của AI nảy sinh từ sự hội tụ của sức mạnh tính toán gia tăng, nguồn dữ liệu khổng lồ và những đột phá thuật toán đáng kể. Ngày nay, khi AI tạo sinh thúc đẩy lĩnh vực này tiến tới những biên giới mới, các bài học từ những mùa đông trước vẫn còn nguyên giá trị: chúng làm nổi bật sự cần thiết phải cân bằng giữa đổi mới và kỳ vọng thực tế, đồng thời giải quyết các thách thức như tính bền vững môi trường, khả năng tiếp cận công bằng và việc duy trì các tiêu chuẩn nghiêm ngặt. Trong thời gian tới, sự hợp tác liên tục giữa các ngành, cùng với việc quản trị có trách nhiệm, sẽ là điều thiết yếu để bảo đảm tiềm năng chuyển đổi của AI mang lại lợi ích cho xã hội một cách rộng rãi và bền vững.
